\documentclass[11pt]{article}
\usepackage[T1]{fontenc}
\usepackage{lmodern,amsmath,amssymb,amsthm,geometry,microtype,xcolor}
\usepackage[colorlinks=true,allcolors=blue!60!black]{hyperref}
\geometry{hmargin=.94in,vmargin=.87in}
\newtheorem{theorem}{Theorem}
\newtheorem{lemma}[theorem]{Lemma}
\newcommand{\R}{\mathbb R}
\newcommand{\Z}{\mathbb Z}
\newcommand{\dimH}{\dim_{\mathrm H}}
\setlength{\emergencystretch}{1em}
\begin{document}
\begin{center}
{\LARGE Hecke limit sets and the Hausdorff dimension\\[3pt]of exceptional square chess billiards}\par\medskip
Blake Stone's mathematics project, developed with OpenAI Codex\\November 27, 2026
\end{center}
\begin{abstract}
The pairs of directions for which square chess billiards has no periodic trajectory form a set of angular measure zero and Hausdorff dimension two. Along each ray of unequal opposite slope magnitudes, their dimension equals that of a Hecke triangle group limit set, with group parameter $(a+b)/\sqrt{ab}$ for slopes $a,-b$. This yields the complete range of slice dimensions and an asymptotic expansion at extreme slope ratios. The proof identifies the chess-billiard exceptional sets, up to countably many points, with the classical limit sets.
\end{abstract}

\section{Direction pairs and Hecke groups}

In the chess billiard of Hanusa and Mahankali~\cite{HM}, a point on the boundary of a square alternates between chords in two prescribed directions. At a supporting corner a zero-length chord counts as one move. Let $\mathcal E$ be the set of nonparallel direction pairs for which there is no periodic orbit, including fixed points as periodic. Use the usual angular coordinates on the two direction circles.

For opposite slopes $a,-b$, with $a,b>0$ and $a\ne b$, put
\[
 l=\frac{\min(a,b)}{a+b},\qquad t=l(1-l),\qquad s=\frac{ab}{a+b},
 \qquad w=\frac{1}{\sqrt t}=\frac{a+b}{\sqrt{ab}}>2.
\]
On either ordering of the slope magnitudes the variables $(l,s)$ form smooth coordinates. Denote by $\mathcal E_l$ the set of $s>0$ producing no periodic trajectory. Write $\Lambda_w$ for the limit set of the Hecke group
\[
 \Gamma_w=\langle T(z)=z+w,\ S(z)=-1/z\rangle
\]
on the boundary of the upper half-plane, and $\delta(w)=\dimH\Lambda_w$.

\begin{theorem}\label{main}
The set $\{s/\sqrt t:s\in\mathcal E_l\}$ and the positive part of $\Lambda_w$ differ by at most countably many points. In particular
\[
 d(l):=\dimH\mathcal E_l=\delta\!\left(\frac{1}{\sqrt{l(1-l)}}\right).
\]
The function $d$ increases continuously and strictly from $1/2$ to $1$ as $l$ ranges from $0$ to $1/2$. With $\gamma$ denoting the Euler--Mascheroni constant, as $t\downarrow0$,
\[
 d(l)=\frac12+\sqrt t+t\log t+2\gamma t
       +O\!\left(t^{3/2}\log^2(1/t)\right).
\]
The set $\mathcal E$ has Hausdorff dimension two, including on every parallelogram after transporting directions by an invertible affine map.
\end{theorem}

The measure-zero assertion is proved in the companion paper~\cite{periodic}. The spectral and asymptotic facts about $\Gamma_w$ are due to Phillips--Sarnak and Soares~\cite{PS,Soares}. The identification of the square's individual rays with these groups is the step that supplies the geometric interpretation and the dimension of the two-parameter set.

\section{Coding the slice}

We use the following square-billiard criterion~\cite[Theorem~1]{periodic}. Start with $x=\operatorname{dist}(s,\Z)$. The parameter belongs to $\mathcal E_l$ precisely if $0<x<l$ at every step of
\[
 x\longmapsto \operatorname{dist}(t/x,\Z).
\]
For context, unfolding across the square's edges produces a monotone trajectory in the unit grid with alternating slopes $a,b$ in the two checkerboard colors. Its vertical drift $D(s)$ is rational exactly when a periodic trajectory exists. Integer translation, reflection and exchange of the axes yield $D(s+j)=D(s)+2j$, $D(-s)=-D(s)$ and $D(t/s)=1/D(s)$. On $[l,1/2]$ the drift equals one, and on $(0,l)$ it lies strictly between zero and one. The algorithm follows from strict descent of the denominator of a rational drift at each nonterminating step.

On $J=[0,l]$ the inverse branches of the interior part of the criterion are the contractions
\[
 f_{j,\epsilon}(x)=\frac{t}{j+\epsilon x},\qquad j\ge1,\quad \epsilon\in\{-1,1\}.
\]
Indeed $t/x>1-l>1/2$ if $0<x<l$, and the nearest integer is positive. The images of $J$ lie in $J$, have disjoint interiors and have derivative magnitude bounded by $l/(1-l)<1$. Infinite digit sequences therefore determine unique points of $J$. Each nonterminating point is obtained in this way. An infinite sequence can encode a terminating point only if one of its tails represents $0$ or $l$. Indeed every tail lies in $J$, and an interior image has the prescribed next remainder by the nearest-integer rule. Such points belong to the countable collection of all finite words applied to $0,l$. The remaining terminal remainders in $(l,1/2]$ are outside the images of the inverse coding.

\begin{lemma}\label{coding}
Up to countably many points, $\Lambda_w$ is the union over $k\in\Z$ of $kw+K$, where $K$ consists of all the limits of infinite compositions of $Q_n(z)=-1/(nw+z)$ with $n\in\Z\setminus\{0\}$. It is symmetric about zero. The image of $K$ under $z\mapsto\sqrt t\,|z|$ equals the set of infinite inverse-branch codes in $J$.
\end{lemma}
\begin{proof}
On the closed unit disk each $Q_n$ takes values in that disk and has Lipschitz constant at most $(w-1)^{-2}<1$, by the difference formula for reciprocals. Thus $K$ is independent of the point from which the compositions are evaluated; using zero shows $K$ is real. Any element of the orbit of $i$ under $\Gamma_w$ can be written
\[
 T^k Q_{n_1}\cdots Q_{n_h}(i).
\]
Indeed $Q_n=ST^n$, one can combine successive powers of $T$ and remove $S^2$, and a trailing $S$ fixes $i$. The limit set is the collection of the boundary accumulation points of this orbit.

For a convergent sequence of orbit points, an unbounded leading translation yields $\infty$ after passing to a subsequence. Otherwise fix $k$ on a subsequence. One may then successively fix each next digit on an infinite subsequence, or choose its absolute value tending to infinity. A sequence of fixed finite depth and fixed digits gives an interior point. A digit tending to infinity maps the entire disk uniformly to zero and yields a translate of a finite word applied to zero. There are only countably many such points. If digits keep stabilizing indefinitely, the common-prefix contraction implies convergence to a point of $kw+K$. Conversely the partial compositions applied to $i$ converge to their real infinite-word limit, which is therefore a limit point. The sign relation $Q_{-n}(-z)=-Q_n(z)$ gives symmetry of $K$.

In signed coordinates $v=\sqrt t\,z$ the map $Q_n$ becomes $v\mapsto -t/(n+v)$. The interval $[-l,l]$ is invariant because $t/(1-l)=l$. Taking absolute values gives $f_{|n|,\epsilon}(|v|)$ with $\epsilon=\operatorname{sgn}(n)\operatorname{sgn}(v)$, with either choice at zero. Conversely, for any folded word $(j_i,\epsilon_i)_{i\ge1}$ choose an initial output sign $\sigma_0$, set $n_i=-\sigma_{i-1}j_i$ and choose the next input sign $\sigma_i=-\sigma_{i-1}\epsilon_i$. This signed word folds to the prescribed one, by the same equality of branches and contraction.
\end{proof}

Since the positive parameter $s$ is first reduced by distance to the integers, its nonterminating set is the positive part of all integer translates of the signed, folded-back central set. The scaling by $1/\sqrt t$ turns those translates into $kw$. Lemma~\ref{coding} proves the countable-difference assertion. The positive part of the limit set has the full Hausdorff dimension by symmetry and translation invariance. Phillips and Sarnak's dimension law for the Hecke groups~\cite{PS} gives strict continuous decrease from one to $1/2$ for $2<w<\infty$. The substitution $w=1/\sqrt t$ in Soares's asymptotic formula~\cite[Theorem~1.9]{Soares} proves the expansion. For example $l=(1-\sqrt5/3)/2$ gives $w=3$ and the rigorous bounds $0.75065<d(l)<0.75322$ from~\cite[Remark~1.12]{Soares}.

\section{Passing from rays to pairs}

Fix any $d<1$ with $d>0$. On an open interval of $l$ below $1/2$, the dimension of each $s$-slice in $(0,1)$ exceeds $d$. Restriction to $(0,1)$ preserves the slice dimension by integer translation. If this subset of the $(l,s)$-plane had dimension less than $1+d$, it would admit, for every $k\ge1$, a countable cover by balls of diameters $r_{k,i}<1/k$ with $\sum_i r_{k,i}^{1+d}<2^{-k}$. At each $l$, intersect these balls with its vertical line. Integrating their summed $d$th powers over $l$ gives at most $\sum_i r_{k,i}^{1+d}$, since each ball projects to an interval of length at most its diameter. Summing in $k$ gives a finite integral. For almost every $l$ the slice covers would then have $d$th powers summing to zero in the limit, contrary to the slice dimension bound on an open interval. The two-parameter dimension is consequently at least $1+d$ for all $d<1$, and so equals two.

The changes between $(l,s)$, the slopes and the two angles are smooth diffeomorphisms on open charts. They preserve Hausdorff dimension by local bi-Lipschitz bounds and countable exhaustion. For any parallelogram the direction map induced by an invertible affine image of the square is likewise a smooth diffeomorphism. This completes the proof of Theorem~\ref{main}.

\small
\begin{thebibliography}{9}
\bibitem{HM} C. R. H. Hanusa and A. V. Mahankali, \emph{A billiards-like dynamical system for attacking chess pieces}, European J. Combin. \textbf{95} (2021), 103341. \href{https://doi.org/10.1016/j.ejc.2021.103341}{doi:10.1016/j.ejc.2021.103341}.
\bibitem{periodic} \emph{Periodic directions in square chess billiards have full measure}, companion proof manuscript, November 2026. Included in this collection as \texttt{manuscript.pdf}.
\bibitem{PS} R. S. Phillips and P. Sarnak, \emph{On the spectrum of the Hecke groups}, Duke Math. J. \textbf{52} (1985), 211--221. \href{https://doi.org/10.1215/S0012-7094-85-05212-3}{doi:10.1215/S0012-7094-85-05212-3}.
\bibitem{Soares} L. Soares, \emph{Hecke triangle groups, transfer operators and Hausdorff dimension}, Ann. Henri Poincar\'e \textbf{23} (2022), 1239--1281. \href{https://doi.org/10.1007/s00023-021-01117-1}{doi:10.1007/s00023-021-01117-1}.
\end{thebibliography}
\end{document}
